Put $\alpha=\angle BAD=\angle ADC$ and $\gamma=\angle ADB$. Then $\angle ABD=\pi-\alpha-\gamma$, $\angle BDC=\alpha-\gamma$, and $\angle DBC=2\gamma$. In particular $0<\gamma<\dfrac{\pi}{2}$.
The sine rule in the two triangles cut by $BD$ gives
\[
\dfrac{\sin\alpha}{\sin\gamma}=\dfrac{BD}{AB}=\dfrac{5}{4},\qquad
\dfrac{\sin(\alpha-\gamma)}{\sin(\alpha+\gamma)}=\dfrac{BC}{BD}=\dfrac{3}{5}.
\]
Expanding the second relation yields $\sin\alpha\cos\gamma=4\cos\alpha\sin\gamma$. Therefore
\[
\sin(\alpha+\gamma)=\sin\alpha\cos\gamma+\cos\alpha\sin\gamma
=\dfrac{5}{4}\sin\alpha\cos\gamma
=\dfrac{25}{16}\sin\gamma\cos\gamma.
\]
Applying the sine rule to $BCD$,
\[
CD=10\,\dfrac{\sin2\gamma}{\sin(\alpha+\gamma)}
=10\cdot\dfrac{2}{\dfrac{25}{16}}=\dfrac{64}{5}.
\]
Thus $\boxed{m+n=69}$.
\begin{asy}
texpreamble("\usepackage{amsmath,amssymb}");
size(240);
defaultpen(fontsize(9pt));
pair B=(0,0),D=(10,0),A=(5.2,2sqrt(231)/5),C=(-174.0/125,-sqrt(36-(174.0/125)^2));draw(A--B--C--D--cycle);draw(B--D,dashed);label("$A$",A,N);label("$B$",B,W);label("$C$",C,SW);label("$D$",D,E);label("$8$",(A+B)/2,NW);label("$10$",(B+D)/2,N);label("$6$",(B+C)/2,W);
\end{asy}
